\documentclass[11pt,a4paper]{article}

\usepackage[T1]{fontenc}
\usepackage[utf8]{inputenc}
\usepackage{amsmath,amssymb}
\usepackage{geometry}
\usepackage{parskip}

\geometry{margin=2.5cm}

\title{A density-characteristic temperature $T_n$ from the logarithmic
  slope of $f$\\ for the \texttt{isc=2} self-collision background}
\author{FP2D\_QLRF\_NL --- design note}
\date{\today}

\begin{document}
\maketitle

This note gives a complete, self-consistent definition of a
density-characteristic temperature $T_n$ that can be used in place of the
energy-weighted $T_\mathrm{eff}$ as the background temperature of the
\texttt{isc=2} self-collision operator, with the velocity convention pinned
down so it is a drop-in replacement.

% ---------------------------------------------------------------------------
\section{Convention (what the operator actually uses)}
% ---------------------------------------------------------------------------

\texttt{cblin} builds its background Maxwellian as the standard Maxwellian
\begin{equation}
  f_M \propto \exp\!\left(-\frac{v^2}{2v_\mathrm{th}^2}\right),
  \qquad
  v_\mathrm{th} = 9.79\times10^3\,\sqrt{T[\mathrm{eV}]/A}\;=\;\sqrt{T/m},
  \qquad m = A\,m_p .
\end{equation}
(In \texttt{cblin}, \texttt{arg = v1/($\sqrt2$\,vth)} and
\texttt{arg2 = v\textasciicircum2/(2\,vth\textasciicircum2)} --- i.e.\ the
Maxwellian at temperature $T$.) A temperature is therefore turned into a
background through $v_\mathrm{th}^2 = T/m$. We define $T_n$ so that it plugs
straight into that same mapping; then it is a drop-in replacement for
$T_\mathrm{eff}$ in the \texttt{isc=2} branch.

% ---------------------------------------------------------------------------
\section{Building block: the local log-slope temperature}
% ---------------------------------------------------------------------------

For that Maxwellian, $\ln f = \mathrm{const} - v^2/2v_\mathrm{th}^2$, so the
radial logarithmic slope is
\begin{equation}
  \frac{\partial \ln f}{\partial v} = -\frac{v}{v_\mathrm{th}^2}
  \quad\Longrightarrow\quad
  v_\mathrm{th}^2(\mathbf v) \equiv \frac{-v}{\partial_v \ln f}
  = \frac{-v^2}{\mathbf v\cdot\nabla_v \ln f},
  \qquad
  T_\mathrm{loc}(\mathbf v) = m\,v_\mathrm{th}^2(\mathbf v).
\end{equation}
This is the pointwise ``temperature seen in the slope of $f$''. It is exact
(constant, equal to $T$) for a Maxwellian and varies across velocity space for
a non-Maxwellian $f$.

% ---------------------------------------------------------------------------
\section{The scalar $T_n$ (density-weighted log-slope)}
% ---------------------------------------------------------------------------

Average the \emph{inverse} local thermal speed, weighted by the distribution
itself (so the dense bulk dominates):
\begin{equation}
  \boxed{\;
  \frac{1}{v_{\mathrm{th},n}^2}
  \;\equiv\;
  \frac{1}{n}\int f\;\frac{-\,\mathbf v\cdot\nabla_v \ln f}{v^2}\,d^3v
  \;=\;
  \frac{1}{n}\int \frac{-\,\mathbf v\cdot\nabla_v f}{v^2}\,d^3v
  \;}, \qquad T_n = m\,v_{\mathrm{th},n}^2 .
\end{equation}
The second form (no $1/f$) is the one to implement. Integrating it by parts,
using $\nabla_v\!\cdot(\mathbf v/v^2) = 1/v^2$ (boundary terms vanish), gives a
remarkably clean closed form:
\begin{equation}
  \boxed{\;
  \frac{1}{v_{\mathrm{th},n}^2}
  = \left\langle \frac{1}{v^2}\right\rangle
  \equiv \frac{1}{n}\int \frac{f}{v^2}\,d^3v
  \quad\Longrightarrow\quad
  T_n = \frac{m}{\langle v^{-2}\rangle}
  \;}
\end{equation}

So the ``log-slope temperature'' \emph{is exactly the harmonic mean of
$v^2$}. Both forms are mathematically identical; the $\langle v^{-2}\rangle$
form is preferable numerically (no derivatives of $f$).

\paragraph{Why this is the right ``density-characteristic'' temperature.}
\begin{itemize}
  \item It reduces to $T$ for a Maxwellian: for
        $f\propto\exp(-v^2/2v_\mathrm{th}^2)$ one finds
        $\langle v^{-2}\rangle = 1/v_\mathrm{th}^2$, hence $T_n = T$.
  \item $\langle v^{-2}\rangle$ is dominated by \emph{small} $v$ (the dense,
        cold core), so $T_n$ tracks the bulk and largely ignores the hot
        perpendicular tail. Contrast this with
        $T_\mathrm{eff}\propto\langle v^2\rangle$, which is tail-weighted. For
        the RF case this gives $T_n\approx$ bulk ($\sim$12\,keV) while
        $T_\mathrm{eff}\approx22$\,keV --- exactly the gap that drives the
        \texttt{isc=2} pathology.
\end{itemize}

% ---------------------------------------------------------------------------
\section{Code recipe (2D cylindrical, existing machinery)}
% ---------------------------------------------------------------------------

\begin{equation}
  \langle v^{-2}\rangle
  = \frac{\displaystyle\sum_{i,j}\frac{f_{ij}}{v_{\perp i}^2+v_{\parallel j}^2}\,J_{ij}}
         {\displaystyle\sum_{i,j} f_{ij}\,J_{ij}},
  \qquad
  T_n[\mathrm{eV}] = \frac{A}{(9.79\times10^3)^2\,\langle v^{-2}\rangle},
\end{equation}
where $J_{ij}$ is the existing \texttt{jacob} ($2\pi v_\perp\times$ grid
weights). Concretely, mirroring \texttt{time\_energy}:

\begin{verbatim}
! density-characteristic (log-slope) temperature; code convention vth^2 = T/m
do iv = 1, nperp
  do ip = 1, npar
    v2 = vperp(iv)**2 + vpar(ip)**2
    fint(iv,ip) = f(iv,ip) / max(v2, v2_floor) * jacob(iv,ip)   ! integrand <1/v^2>
  end do
end do
call ncint_2d(fint, mom_invv2)          ! = n * <1/v^2>
inv_v2_avg = mom_invv2 / dens           ! <1/v^2>
Tn_eV = aa / (9.79d3**2 * inv_v2_avg)   ! drop-in for ta_eV/1e3 in the isc=2 branch
\end{verbatim}

\begin{itemize}
  \item \textbf{Origin guard:} the integrand $\propto f/v^2$ has only a mild,
        integrable singularity at $\mathbf v = 0$; since
        \texttt{vperp(1)}~$\approx 394$~m\,s$^{-1} > 0$ there is no exact zero,
        but keep a small \texttt{v2\_floor} (e.g.\ \texttt{vperp(1)**2}) for
        safety.
  \item To use it, replace the \texttt{isc=2} line
        \texttt{vteff\_t = 9.79d3*SQRT(ta\_eV/aa)} by
        \texttt{vteff\_t = 9.79d3*SQRT(Tn\_eV/aa)} (and use \texttt{Tn\_eV} for
        the self-collision Coulomb log). Keep writing both
        \texttt{Teff\_vs\_time} and a new \texttt{Tn\_vs\_time} so that the gap
        $T_\mathrm{eff}/T_n$ is visible.
\end{itemize}

% ---------------------------------------------------------------------------
\section*{Remark: a related convention point}
% ---------------------------------------------------------------------------

The diagnostic SC-Maxwellian output (\texttt{fsc\_maxw} in
\texttt{TimeFP\_7pt.f90}, and \texttt{stix\_maxw\_plot.py}) originally used
$\exp(-v^2/v_\mathrm{th}^2)$ with prefactor $1/(\pi^{3/2}v_\mathrm{th}^3)$ ---
too narrow by $\sqrt2$ and mis-normalised relative to the
$\exp(-v^2/2v_\mathrm{th}^2)$, $1/((2\pi)^{3/2}v_\mathrm{th}^3)$ background that
\texttt{cblin} actually uses. This has since been corrected so the diagnostic
matches the operator. It does not affect the $T_n$ definition above, which is
stated directly in the \texttt{cblin} convention.

\end{document}
